\section{Step 48 Observe the first reduction tile before the second}\label{matmul-s48}

Use the completed tiled program for the default shape and stop at \texttt{block\_reduction} on its second visit, after L becomes two. The first visit merely selects the initial segment. At the second, C contains the first two terms of every output, while the last product has not yet been added. Continue to \texttt{matmul\_done} and compare the final exact values.

Use recipe \texttt{s48-1}. Pause at the second visit to \texttt{block\_reduction}; this is an internal observation before the restoring exit.

\subsection{Expected observation}
Before the second segment, C is 25,28,73,82; at return it is 58,64,139,154.

\Needspace{10\baselineskip}
\subsection{Check your understanding}
Which observation proves that the second tile must retain prior state?

\begin{itemize}
\item \textbf{A} The output allocation is mapped.
\item \textbf{B} The descriptor tile edge is positive.
\item \textbf{C} C already contains valid first-segment contributions.
\end{itemize}
